\documentclass[a4paper,14pt]{extarticle}
\usepackage{cmap}
\usepackage[T2A]{fontenc}
\usepackage[utf8]{inputenc}
\usepackage[russian]{babel}
\usepackage[left=30mm,right=15mm,top=20mm,bottom=23mm]{geometry}
\usepackage{indentfirst}
\usepackage{setspace}
\usepackage{booktabs,tabularx}
\usepackage{caption,listings}
\usepackage{tikz}
\usepackage[hidelinks,unicode]{hyperref}

\setstretch{1.30}
\setlength{\parindent}{1.25cm}
\setlength{\emergencystretch}{1em}
\pagestyle{plain}
\captionsetup{font=small,labelsep=endash}
\captionsetup[table]{position=top}
\renewcommand{\lstlistingname}{Листинг}
\lstset{
  basicstyle=\ttfamily\small,
  numbers=left,
  numberstyle=\scriptsize\color{black!55},
  xleftmargin=1.5em,
  frame=single,
  rulecolor=\color{black!30},
  captionpos=b,
  columns=fullflexible,
  keepspaces=true,
  breaklines=true,
  showstringspaces=false
}
\hypersetup{
  pdftitle={Лабораторная работа № 3. Маршрутизация в GNS3},
  pdfauthor={Храмов В. Д.},
  pdfsubject={Сети ЭВМ и телекоммуникации компьютерных сетей. Вариант 22}
}

\begin{document}

\begin{titlepage}
  \thispagestyle{empty}
  \begin{center}
    \begin{singlespace}
      МИНИСТЕРСТВО НАУКИ И ВЫСШЕГО ОБРАЗОВАНИЯ\\
      РОССИЙСКОЙ ФЕДЕРАЦИИ\par\medskip
      Санкт-Петербургский политехнический университет Петра Великого\par\medskip
      Институт компьютерных наук и кибербезопасности\\
      Высшая школа технологий искусственного интеллекта\\
      Направление 02.03.01 «Математика и компьютерные науки»
    \end{singlespace}

    \vfill
    {\Large\bfseries Отчёт о выполнении\\лабораторной работы №~3\par}
    \medskip
    {\large Динамическая маршрутизация в GNS3\par}
    \medskip
    Дисциплина «Сети ЭВМ и телекоммуникации компьютерных сетей»\par
    Вариант 22\par
  \end{center}

  \vfill
  \noindent Выполнил студент группы 5130201/40101\hfill
  \underline{\hspace{2.3cm}}\enspace Храмов В.~Д.\par\medskip
  \noindent Проверил, к.\,т.\,н.\hfill
  \underline{\hspace{2.3cm}}\enspace Мулюха В.~А.\par\medskip
  \begin{flushright}
    «\underline{\hspace{0.9cm}}» \underline{\hspace{2.6cm}} 2026 г.
  \end{flushright}

  \vfill
  \begin{center}
    Санкт-Петербург\\2026
  \end{center}
\end{titlepage}

\setcounter{page}{2}
\tableofcontents
\clearpage

\section{Цель и постановка задачи}

Цель работы --- построить сеть по схеме варианта~22 в GNS3, разработать план
IP-адресации с использованием как минимум трёх разных масок и проверить
связность конечных и промежуточных узлов. При отключении одной из ветвей
трафик должен проходить по оставшемуся пути после автоматического
пересчёта маршрутов.

На рисунке~\ref{fig:topology} показана реализованная топология. Хосты S1--S3
подключены к отдельным интерфейсам R1, а хосты K1--K3 подключены к R3
через коммутатор SW. Между R1 и R3 построено два пути.

\begin{figure}[htbp]
  \centering
  \begin{tikzpicture}[
    x=1.37cm,y=1.00cm,
    every node/.style={font=\small},
    router/.style={circle,draw,thick,minimum size=9mm,inner sep=0pt},
    host/.style={rectangle,draw,minimum width=10mm,minimum height=7mm},
    switch/.style={rectangle,draw,very thick,minimum width=11mm,minimum height=7mm},
    local/.style={draw=black!60},
    primary/.style={draw=blue!60!black,very thick},
    backup/.style={draw=blue!60!black,thick,dashed}
  ]
    \node[router] (r1) at (1.35,0) {R1};
    \node[router] (r2) at (3.10,0) {R2};
    \node[router] (r3) at (4.85,0) {R3};
    \node[router] (r4) at (1.35,1.65) {R4};
    \node[router] (r5) at (3.10,1.65) {R5};
    \node[router] (r6) at (4.85,1.65) {R6};
    \node[host] (s1) at (0,1) {S1};
    \node[host] (s2) at (0,0) {S2};
    \node[host] (s3) at (0,-1) {S3};
    \node[switch] (sw) at (6.30,0) {SW};
    \node[host] (k1) at (7.65,1) {K1};
    \node[host] (k2) at (7.65,0) {K2};
    \node[host] (k3) at (7.65,-1) {K3};

    \draw[primary] (r1) -- (r2) -- (r3);
    \draw[backup] (r1) -- (r4) -- (r5) -- (r6) -- (r3);
    \draw[local] (s1) -- (r1) (s2) -- (r1) (s3) -- (r1);
    \draw[local] (r3) -- (sw) (sw) -- (k1) (sw) -- (k2) (sw) -- (k3);
  \end{tikzpicture}
  \caption{Схема сети варианта 22. Пунктиром показана резервная ветвь}
  \label{fig:topology}
\end{figure}

\section{Теоретические сведения}

OSPFv2 распространяет сведения о состоянии каналов внутри области и
вычисляет пути с минимальной суммарной стоимостью исходящих
интерфейсов. После изменения топологии OSPF обновляет базу данных
состояния каналов и пересчитывает таблицу маршрутизации~\cite{rfc2328}.
На каналах между маршрутизаторами используется режим
\texttt{point-to-point}. Интерфейсы, ведущие к хостам, настроены как
пассивные. Их сети объявляются в OSPF, а соседства на этих интерфейсах
не формируются~\cite{cisco-ospf}.

\clearpage
\section{Реализация}

\subsection{Оборудование и план адресации}

Использованы шесть Cisco~3745 с IOS~12.4(25) в Dynamips, шесть узлов
VPCS и встроенный Ethernet Switch GNS3. Cisco~3745 располагает двумя
встроенными Fast Ethernet и поддерживает дополнительные сетевые
модули~\cite{gns3-ios}. На R1 установлены три модуля NM-1FE-TX для хостов
S1--S3. На R3 один такой модуль соединяет маршрутизатор с SW. Порт
Ethernet0 коммутатора подключён к R3, Ethernet1--Ethernet3 --- к K1--K3
соответственно. Все четыре порта работают в режиме access в VLAN~1.

План адресации использует пять масок. В таблице~\ref{tab:lan} указаны
четыре сети конечных узлов. Шлюз каждого хоста --- адрес интерфейса
маршрутизатора в его локальной сети.

\begin{table}[htbp]
  \centering\small
  \caption{Локальные сети и адреса конечных узлов}
  \label{tab:lan}
  \begin{tabularx}{\linewidth}{@{}llXX@{}}
    \toprule
    Подсеть & Интерфейс & Шлюз & Адреса хостов \\
    \midrule
    192.168.22.0/24  & R1 Fa1/0 & 192.168.22.1  & S1 192.168.22.2 \\
    192.168.23.0/26  & R1 Fa2/0 & 192.168.23.1  & S2 192.168.23.2 \\
    192.168.23.64/27 & R1 Fa3/0 & 192.168.23.65 & S3 192.168.23.66 \\
    192.168.24.0/28  & R3 Fa1/0 & 192.168.24.1  & K1 192.168.24.2 \\
                     &          &               & K2 192.168.24.3 \\
                     &          &               & K3 192.168.24.4 \\
    \bottomrule
  \end{tabularx}
\end{table}

Межмаршрутизаторные каналы перечислены в таблице~\ref{tab:transit}.
Маска /30 выделяет два адреса узлов на каждом канале.

\begin{table}[htbp]
  \centering\small
  \caption{Адреса межмаршрутизаторных каналов}
  \label{tab:transit}
  \begin{tabularx}{\linewidth}{@{}lXX@{}}
    \toprule
    Подсеть & Первый интерфейс & Второй интерфейс \\
    \midrule
    10.22.0.0/30  & R1 Fa0/0 --- 10.22.0.1  & R2 Fa0/0 --- 10.22.0.2 \\
    10.22.0.4/30  & R2 Fa0/1 --- 10.22.0.5  & R3 Fa0/0 --- 10.22.0.6 \\
    10.22.0.8/30  & R1 Fa0/1 --- 10.22.0.9  & R4 Fa0/0 --- 10.22.0.10 \\
    10.22.0.12/30 & R4 Fa0/1 --- 10.22.0.13 & R5 Fa0/0 --- 10.22.0.14 \\
    10.22.0.16/30 & R5 Fa0/1 --- 10.22.0.17 & R6 Fa0/0 --- 10.22.0.18 \\
    10.22.0.20/30 & R6 Fa0/1 --- 10.22.0.21 & R3 Fa0/1 --- 10.22.0.22 \\
    \bottomrule
  \end{tabularx}
\end{table}

\clearpage
\subsection{Настройка OSPF}

На R1--R6 запущен процесс OSPF~1 в области~0. Идентификатор
маршрутизатора R\(i\) задан в виде \(i.i.i.i\), где \(i\) --- номер
маршрутизатора от 1 до 6. Команда \texttt{passive-interface default}
отключает формирование соседств на всех интерфейсах по умолчанию.
Соседства разрешены только на транзитных каналах. На них интервал
Hello установлен равным 1~с, а Dead --- 4~с. Значения совпадают на обоих
концах каждого канала.

Стоимость интерфейса явно задана равной 10. До сети 192.168.24.0/28
из R1 стоимость основного пути составляет \(2\cdot10+10=30\), а
резервного --- \(4\cdot10+10=50\). Последнее слагаемое в обоих
случаях соответствует интерфейсу R3 в сети K1--K3. При доступности
обеих ветвей OSPF выбирает путь через R2.

В листинге~\ref{lst:r1} приведён фрагмент стартовой конфигурации R1.
Команды \texttt{network} используют точное совпадение с адресом
интерфейса. OSPF объявляет соответствующую подключённую подсеть с
её фактической маской.

{\singlespacing
\begin{lstlisting}[caption={Фрагмент конфигурации OSPF на R1},label={lst:r1}]
router ospf 1
 router-id 1.1.1.1
 log-adjacency-changes
 passive-interface default
 no passive-interface FastEthernet0/0
 network 10.22.0.1 0.0.0.0 area 0
 no passive-interface FastEthernet0/1
 network 10.22.0.9 0.0.0.0 area 0
 network 192.168.22.1 0.0.0.0 area 0
 network 192.168.23.1 0.0.0.0 area 0
 network 192.168.23.65 0.0.0.0 area 0
!
interface FastEthernet0/0
 ip address 10.22.0.1 255.255.255.252
 ip ospf cost 10
 ip ospf network point-to-point
 ip ospf hello-interval 1
 ip ospf dead-interval 4
 duplex auto
 speed auto
 no shutdown
\end{lstlisting}
}

\clearpage
\section{Проверка работы сети}

\subsection{Штатный режим}

Для проверки использованы GNS3 2.2.61, Dynamips 0.2.23 и
VPCS 0.8.3. Эксперимент проведён 22~ноября 2026~г. Архив проекта
включает конфигурации узлов и образ IOS из материалов к заданию.
Проверки выполнены после загрузки
узлов и формирования соседств OSPF. В каждой серии \texttt{ping} отправлялось по пять запросов ICMP Echo Request.

Команда \texttt{show ip ospf neighbor} показала по два соседства
в состоянии \texttt{FULL} на каждом маршрутизаторе. В
листинге~\ref{lst:neighbors} приведён вывод на R1. Сосед R2 находится
на основной ветви, а R4 --- на резервной.

{\singlespacing
\begin{lstlisting}[caption={Соседства OSPF на R1},label={lst:neighbors},basicstyle=\ttfamily\footnotesize]
R1#show ip ospf neighbor
Neighbor ID     Pri   State           Dead Time   Address         Interface
4.4.4.4           0   FULL/  -        00:00:03    10.22.0.10      FastEthernet0/1
2.2.2.2           0   FULL/  -        00:00:03    10.22.0.2       FastEthernet0/0
\end{lstlisting}
}

В таблице~\ref{tab:ping} приведены проверки всех шести хостов и
промежуточных маршрутизаторов R2, R4, R5 и R6. Через R1 и R3
проходил сквозной трафик, что подтвердили трассировка и таблицы
маршрутизации. В таблице приведены итоговые серии. При первом обмене S1--K1 получено 4/5 ответов, при повторе --- 5/5.

\begin{table}[htbp]
  \centering\small
  \caption{Результаты ping в штатном режиме}
  \label{tab:ping}
  \begin{tabularx}{\linewidth}{@{}lXr@{}}
    \toprule
    Источник & Получатель & Ответы / запросы \\
    \midrule
    S1 & K1 (192.168.24.2) & 5/5 \\
    S2 & K2 (192.168.24.3) & 5/5 \\
    S3 & K3 (192.168.24.4) & 5/5 \\
    K1 & S1 (192.168.22.2) & 5/5 \\
    R2 & R3 Fa0/0 (10.22.0.6) & 5/5 \\
    R4 & R5 Fa0/0 (10.22.0.14) & 5/5 \\
    R5 & R6 Fa0/0 (10.22.0.18) & 5/5 \\
    R6 & R3 Fa0/1 (10.22.0.22) & 5/5 \\
    \bottomrule
  \end{tabularx}
\end{table}

\clearpage
\subsection{Отключение ветвей}

При остановке R5 соседства с ним исчезли на R4 и R6. Путь через R2
сохранился. После запуска R5 и восстановления соседств был остановлен
R2. OSPF автоматически перенаправил трафик с R1 через R4, а обратный
трафик с R3 --- через R6. После запуска R2 вернулся путь с меньшей
стоимостью.

\begin{table}[htbp]
  \centering\small
  \caption{Маршрут R1 к сети K1--K3 и доступность K1 с S1}
  \label{tab:failover}
  \begin{tabularx}{\linewidth}{@{}Xlrr@{}}
    \toprule
    Состояние & Следующий узел & Метрика & Ответы / запросы \\
    \midrule
    Все узлы запущены   & R2 (10.22.0.2)  & 30 & 5/5 \\
    R5 остановлен       & R2 (10.22.0.2)  & 30 & 5/5 \\
    R2 остановлен       & R4 (10.22.0.10) & 50 & 5/5 \\
    R2 снова запущен    & R2 (10.22.0.2)  & 30 & 5/5 \\
    \bottomrule
  \end{tabularx}
\end{table}

В таблице~\ref{tab:failover} приведены результаты после сходимости
OSPF. При остановленном R2 доля ответов составила 5/5 для S2--K2,
5/5 для S3--K3 и 5/5 для K1--S1. После запуска
R2 первая серия S1--K1 дала 2/5 ответов. В повторной серии
получены все пять ответов. Длительность перерыва отдельно не измерялась.

В листинге~\ref{lst:trace} приведены строки таблицы R1 и переходы
трассировки S1--K1. При остановке R2 число маршрутизаторов на пути
увеличилось с трёх до пяти.

{\singlespacing
\begin{lstlisting}[caption={Маршрут R1 и переходы трассировки до и после остановки R2},label={lst:trace},basicstyle=\ttfamily\footnotesize]
R1#show ip route 192.168.24.0
O       192.168.24.0 [110/30] via 10.22.0.2, 00:00:00, FastEthernet0/0
S1> trace 192.168.24.2
1  192.168.22.1
2  10.22.0.2
3  10.22.0.6
4  192.168.24.2

R1#show ip route 192.168.24.0
O       192.168.24.0 [110/50] via 10.22.0.10, 00:00:00, FastEthernet0/1
S1> trace 192.168.24.2
1  192.168.22.1
2  10.22.0.10
3  10.22.0.14
4  10.22.0.18
5  10.22.0.22
6  192.168.24.2
\end{lstlisting}
}

\clearpage
\subsection{Порядок демонстрации}

\begin{enumerate}
  \item Открыть меню \texttt{File}, выбрать \texttt{Import portable project}
    и указать файл \path{Lab3.gns3project}. Запустить все узлы.
  \item На R1 и R3 выполнить \texttt{show ip ospf neighbor}. На R1
    вывести \texttt{show ip route 192.168.24.0}, на R3 ---
    \texttt{show ip route 192.168.22.0}. На S1 выполнить
    \texttt{ping 192.168.24.2 -c 5} и \texttt{trace 192.168.24.2}.
  \item Остановить R5 командой \texttt{Stop} из контекстного меню
    узла. Повторить \texttt{ping} и показать, что путь идёт через R2.
  \item Запустить R5 и дождаться двух соседств \texttt{FULL} на нём.
    Остановить R2, проверить смену маршрутов и повторить \texttt{ping}
    и \texttt{trace} с S1.
  \item Запустить R2 и показать возврат маршрута через него.
\end{enumerate}

\section{Заключение}

В GNS3 реализована схема варианта~22 с шестью маршрутизаторами,
одним коммутатором и шестью конечными узлами. В плане адресации
использовано пять длин префикса. Проверки подтвердили доступность
хостов и промежуточных узлов. Остановка R5 оставила доступным путь
через R2, а остановка R2 вызвала автоматическое переключение на
R4--R5--R6. После восстановления R2 снова был выбран
маршрут с меньшей метрикой.

\phantomsection
\addcontentsline{toc}{section}{Список литературы}
{\singlespacing\small
\begin{thebibliography}{9}
  \setlength{\itemsep}{0.25em}
  \bibitem{rfc2328}
    Moy J. \href{https://www.rfc-editor.org/rfc/rfc2328}{OSPF Version 2. RFC 2328.} 1998.
  \bibitem{cisco-ospf}
    Cisco. \href{https://www.cisco.com/c/en/us/td/docs/ios-xml/ios/iproute_ospf/command/iro-cr-book/ospf-a1.html}{OSPF Commands. Cisco IOS IP Routing Command Reference}.
  \bibitem{gns3-ios}
    GNS3. \href{https://docs.gns3.com/docs/emulators/cisco-ios-images-for-dynamips}{Cisco IOS images for Dynamips}.
\end{thebibliography}
}

\end{document}
